% =====================================================================
\section{引言}
% =====================================================================
本课程设计的题目来自一场真实的线下游戏。上海大学逐浪徒步社团在宝山校区 A–G 楼一带（含 AJ–GJ 楼）组织了一次 30 人的猫捉老鼠：开局 10 只猫、20 只老鼠，碰一下就算抓到，没有安全区；被抓的老鼠立刻变成猫，所有人都知道谁变了。游戏分上下两轮，每轮 30 分钟，开局时所有人在 C、D、CJ、DJ 四栋楼之间的空地集合。活动范围只限楼外的平地，不能进教学楼，因为当时楼里还在上课，而且楼里地形复杂，容易受伤。抓得最多的猫和活得最久的老鼠各有奖励。笔者在第一轮担任老鼠。

玩下来有三条直观经验。第一，活得最久的老鼠都是\textbf{躲}出来的：有人躲在灌木丛里，有人躲在空调外机和教学楼外墙之间的夹缝里。躲在夹缝里的那位一直到最后都没有被抓到，成了这一轮的“鼠王”。第二，猫有时几只一起围一只老鼠，有时一只追一只，也有一只猫追好几只、结果被分头甩开。第三，晚上有夜跑的路人，老鼠常把他们错认成猫，吓一跳。

这些经验背后有几个博弈问题。躲和跑哪个更好，取决于什么条件？一个只想多活一会儿的老鼠，会不会自己学会躲？“抓得最多的猫有奖”是一种个人激励，猫之间既要合作围捕，又在争抢功劳，奖励规则会不会改变猫之间的合作？还有，游戏有上下两轮：第一轮跑的人如果看到躲的人活得更久，第二轮会不会改成躲？反过来，会不会有人从躲改成跑？据此，本文提出五个研究问题：

\begin{enumerate}
  \item \textbf{机制（RQ1）}：在这片校园里，“躲比跑活得久”还成立吗？靠的是“站着不动不显眼”，还是“藏身处不显眼”？
  \item \textbf{最好的藏身点（RQ2）}：灌木丛和夹缝哪个更好，和“鼠王躲在夹缝里”的实战经验一致吗？
  \item \textbf{上下两轮之间的转化（RQ3）}：看到第一轮的结果后，老鼠会从跑转向躲吗？会出现从躲转向跑的反方向转化吗？
  \item \textbf{老鼠的学习（RQ4）}：只奖励“活得久”、和猫一起学习的 AI 老鼠，会不会自己学会躲？如果藏身处不再有隐蔽优势，它会不会改成跑？它学到的本领换个对手还管用吗？
  \item \textbf{猫的学习（RQ5）}：不把围捕写进规则，猫能不能自己学会？个人奖励和团队奖励下，猫的强弱和围捕有什么不同？
\end{enumerate}

% =====================================================================
\section{相关工作}
% =====================================================================
表~\ref{tab:related} 先给出相关工作与本文的对应关系，下面再逐类说明。

\begin{table}[htbp]
  \centering\small
  \caption{相关工作与本文的对应关系}
  \label{tab:related}
  \renewcommand{\arraystretch}{1.25}
  \begin{tabularx}{\textwidth}{@{}p{3.1cm}X X@{}}
    \toprule
    文献 & 研究的问题 & 本文借鉴或对照的地方\\
    \midrule
    搜索博弈\cite{alpern2003,chung2011} & 藏的一方选位置，找的一方选搜索路线 & 发现模型（第~\ref{sec:detect} 节）；猫的巡逻和搜索策略\\
    多智能体捉迷藏\cite{baker2020} & 藏者和找者一起训练，会自己学出新策略 & 猫和老鼠一起训练，看老鼠会不会自己学会躲（RQ4）\\
    PPO 与 GAE\cite{schulman2017ppo,schulman2016gae,yu2022mappo} & 稳定的策略梯度训练方法 & 本文的训练算法（第~\ref{sec:obs} 节）\\
    追捕中的奖励结构\cite{grupen2022,leibo2017} & 个人奖励和共享奖励下，追捕者会不会合作 & 猫的个人奖励与团队奖励对比（RQ5）\\
    演化博弈与模仿\cite{taylor1978,szabo1998,traulsen2006} & 人们看到别人的收益后怎样改策略 & 上下两轮之间的策略转化（RQ3）\\
    传染过程\cite{kermack1927} & 接触后“被感染”的人数怎样增长 & 被抓的老鼠变成猫，猫越来越多\\
    \bottomrule
  \end{tabularx}
\end{table}

\subsection{捉迷藏与搜索博弈}
捉迷藏在博弈论里叫搜索博弈：藏的一方选一个位置，找的一方选一条搜索路线，双方都希望对方猜不中\cite{alpern2003}。机器人领域把它和追捕问题合在一起研究，并指出视线被障碍挡住时，追捕会先变成“找人”，找到以后才是“追人”\cite{chung2011}。本文的游戏正是如此：老鼠被发现之前，是猫在找；被发现之后，才是猫在追。“躲”的作用是让老鼠尽量停留在“找”的阶段。本文用一个发现模型（第~\ref{sec:detect} 节）来描述“找”的难度，用猫的巡逻和搜索策略来描述“找”的方式。

\subsection{多智能体强化学习中的捉迷藏}
与本文最接近的是 OpenAI 的捉迷藏实验\cite{baker2020}：藏者和找者用强化学习一起训练，双方先后学出了堵门、搬箱子等策略，而这些策略都没有写进规则。本文借用这种“双方一起学”的做法，看只奖励“活得久”的老鼠会不会自己学会躲。训练算法采用近端策略优化（PPO）\cite{schulman2017ppo} 和广义优势估计（GAE）\cite{schulman2016gae}；同一阵营的个体共用一套网络参数，这种做法在多智能体任务中很常用\cite{yu2022mappo}。需要说明的是，本文的价值网络只用每个人自己的观测，属于“按阵营共享参数的独立 PPO”，不是集中式训练。

\subsection{奖励结构与合作}
Grupen 等人研究了多只追捕者抓一只逃跑者的任务\cite{grupen2022}：追捕者只拿自己的奖励时，倾向于各自贪心地追；改成共享奖励后，更容易形成包抄。他们也指出，共享奖励下抓捕数会集中在少数个体身上，这可能是分工（有人负责堵、有人负责抓），不一定是偷懒。Leibo 等人则说明，个体利益与集体利益的冲突会出现在多步决策中\cite{leibo2017}。本文的猫既要合作围捕，又在比谁抓得多，正好可以检验奖励方式对合作的影响。

\subsection{演化博弈与模仿}
演化博弈研究的是策略怎样在人群中传播：收益高的策略会被更多人采用，这一过程可以用复制者动态描述\cite{taylor1978}。在个体层面，常用的模型是“两两比较”：每个人看一个同伴，同伴收益比自己高时，以一定概率改用同伴的策略，这个概率用费米函数刻画，下文称为费米模仿规则\cite{szabo1998,traulsen2006}。本文用它来回答 RQ3：第一轮结束后，看到别人收益的老鼠会怎样改策略。

\subsection{传染式的转化}
被抓的老鼠变成猫，类似传染病里“易感者接触感染者后被感染”的 SI 过程\cite{kermack1927}：猫越多，剩下的老鼠越危险。这里只是类比，本文的接触过程由移动、视线和策略决定，并没有用流行病方程建模。

% =====================================================================
\section{算法设计}
% =====================================================================
\subsection{环境：基于 OpenStreetMap 的 A–G 楼地图}
地图来自 OpenStreetMap（OSM）\cite{osm}。OSM 是二维地图，记录的是建筑轮廓、道路、林地等平面要素，没有地形起伏，也没有完整的建筑高度，因此本文把场地当作平面，把楼当作完全挡住视线的墙。笔者用 Overpass 接口下载了宝山校区的全部要素（共 19655 个），截取 A–GJ 楼外扩 25 m 的区域，共 418 m × 526 m，按 0.5 m 一格栅格化成 836 × 1053 的格子图（图~\ref{fig:env}a）。

每个格子属于五类之一：平地（路、广场、草坪，共 15.4 万 m²）可以正常行走；教学楼（4.9 万 m²）不能进入，并且挡住视线；灌木丛（1.5 万 m²）取自 OSM 的林地和行道树，可以进入，但只能以 0.8 m/s 慢慢挪；空调外机夹缝（391 m²）；其余是场地外，不能去。下载的 OSM 数据中没有任何与空调外机有关的标签，也就是说 OSM 并没有采集空调外机的位置，因此本文只能按固定随机种子沿 A–GJ 楼外墙放置 235 台空调外机，平均每 14 m 一台，每台和外墙之间留出 2 m 长、1 m 深的夹缝。这一部分是合成的，不代表实际位置。

开局集合点取 C、D、CJ、DJ 四栋楼中心的平均位置附近最近的空地，与实际游戏一致；所有人在集合点周围 30 m × 30 m 范围内的空地上随机出生。寻路方面，笔者在 4 m 的粗格上（11269 个结点）用 Dijkstra 算法预先算好任意两点之间的下一跳。智能体每一步先检查到目标的直线是否通畅，通畅就直走，否则沿最短路绕楼走。
